{\centering \section*{Lời cảm ơn}}

Qua thời gian học tập và rèn luyện tại
\textit{Đại học Bách khoa Hà Nội},
đến nay em đã hoàn thành đồ án tốt nghiệp.
Trong thời gian học tập để có được kết quả hiện tại
em xin chân thành cảm ơn tới các thầy cô giảng viên trong \textit{Khoa Toán Tin}
đã giảng dạy, quan tâm và tạo điều kiện thuận lợi để em học tập rèn luyện trong suốt thời gian qua,
giúp em trang bị những kiến thức, kỹ năng cần thiết
để em có thể hoàn thành tốt nhất báo cáo này
với
đề tài:
\textit{Xây dựng ứng dụng hỗ trợ tư vấn pháp luật sử dụng công nghệ trí tuệ nhân tạo}.

Đồng thời, em xin gửi lời cảm ơn chân thành và sâu sắc đến \textit{Thầy TS. Trần Anh Tú},
người thầy đã tận tình hỗ trợ,
giúp đỡ và hướng dẫn em suốt thời gian thực hiện đồ án.
Những kiến thức và kinh nghiệm mà em đã tiếp thu được trong quá trình này
sẽ đóng góp quan trọng vào sự phát triển và thành công của em trong tương lai.
Em xin gửi lời chúc sức khỏe tốt nhất đến thầy,
hy vọng thầy luôn dồi dào sức khỏe, đam mê và nhiệt huyết trong công việc giảng dạy.

Em cũng xin gửi lời cảm ơn tới bố mẹ
đã dành nhiều thời gian chăm lo, động viên
và giúp đỡ em để em đạt được thành công trong học tập.

Mặc dù đã cố gắng nỗ lực học hỏi để hoàn thành đồ án,
nhưng thời gian thực hiện đồ án có hạn,
kiến thức của em còn hạn chế.
Do vậy,
không tránh khỏi những thiếu sót,
em rất mong nhận được những ý kiến đóng góp quý báu của thầy cô
để em có thể cải thiện một cách tốt nhất.

\vspace{0.5cm} \emph{Em xin chân thành cảm ơn!}

%%%%%%%%%%%%%%%%%%%%%%%%%%%%%%%%%%%%%%%%%%%%%%%%%%%%%%%
%%%%%%%%%%%%%%%%%%%%%%%%%%%%%%%%%%%%%%%%%%%%%%%%%%%%%%%
%%%%%%%%%%%%%%%%%%%%%%%%%%%%%%%%%%%%%%%%%%%%%%%%%%%%%%%

{\centering \section*{Tóm tắt nội dung đồ án}}

% {
% \color{RedHUST}

% Tóm tắt nội dung của đồ án tốt nghiệp trong khoảng tối đa 300 chữ.
% Phần tóm tắt cần nêu được các ý:
% vấn đề cần thực hiện;
% phương pháp thực hiện;
% công cụ sử dụng (phần mềm, phần cứng…);
% kết quả của đồ án có phù hợp với các vấn đề đã đặt ra hay không;
% tính thực tế của đồ án,
% định hướng phát triển mở rộng của đồ án (nếu có);
% các kiến thức và kỹ năng mà sinh viên đã đạt được.

% }

Đồ án
được thực hiện nhằm giải quyết nhu cầu tra cứu
và giải đáp thông tin pháp luật Việt Nam
(bao gồm Văn bản pháp luật và Pháp Điển) một cách chính xác, nhanh chóng và tự động hóa.

Để tối ưu năng lực giải đáp,
đồ án triển khai hệ thống
Agentic RAG nâng cao sử dụng LangGraph.

Về mặt kiến trúc, hệ thống được thiết kế theo kiến trúc vi dịch vụ (Microservice Architecture),
ứng dụng phương pháp Thiết kế hướng miền (Domain-Driven Design - DDD)
và các dịch vụ được điều phối qua Kong API Gateway.

Để đảm bảo tính ổn định,
hiệu năng cao,
tận dụng điểm mạnh tài nguyên,
thế mạnh framework ngôn ngữ lập trình
và khả năng mở rộng,
các dịch vụ giao tiếp đa dạng từ gRPC cho trò chuyện trực tiếp,
giao tiếp bất đồng bộ giữa các dịch vụ NestJS và Python
được xử lý thông qua RabbitMQ.
Luồng nghiệp vụ
quy trình thanh toán thuê bao VIP
qua các cổng nội địa (VNPay, MoMo, ZaloPay, SePay),
được đồng bộ
sự kiện
bằng nền tảng Kafka.
Toàn bộ vòng đời ứng dụng
được tự động hóa thông qua
CI/CD trên GitHub Actions,
đóng gói bằng Docker và triển khai vận hành trên
cụm Kubernetes thông qua ArgoCD.

Đồ án đã vận hành thành công một trợ lý pháp luật,
đưa ra các tư vấn pháp lý với các bước suy luận logic
và trích dẫn nguồn văn bản rõ ràng.
Bên cạnh đó,
ứng dụng còn mang đến trải nghiệm đàm thoại
thời gian thực bằng cách tích hợp LiveKit
cùng các mô hình xử lý giọng nói tiên tiến như Whisper (Speech-to-Text), Edge, Piper và ElevenLabs (Text-to-Speech).

\begin{flushright}
\begin{tabular}{c}
% Hà Nội, \today \\
\textit{Hà Nội, ngày \dots tháng \dots năm \dots} \\
Sinh viên thực hiện \\
\textit{(Ký và ghi rõ họ tên)} \\
% Nghĩa \\
% Vũ Văn Nghĩa
\end{tabular}
\end{flushright}
